\section{Experimental Protocol}
\label{sec:experiments}
\label{sec:empirical_evaluation}

\subsection{Corpora, Targets and Attackers}
\label{sec:setup}

Single-testbed evaluations have limited external validity~\cite{cantone2024crossdataset}, and attack evaluations that ignore realism overstate the risk~\cite{apruzzese2022modeling, apruzzese2023real}, so we use three widely studied corpora in full: CSE-CIC-IDS2018~\cite{sharafaldin2018toward} (nine daily files, seven capped at 1,048,575 rows in the distribution), CIC-IDS2017~\cite{sharafaldin2018toward} (five daily files) and UNSW-NB15~\cite{moustafa2015unsw} (the official training and testing partition, 257,673 of the 2.54 million records). Every record is mapped onto the 12-feature schema of Section~\ref{sec:formulation}, invalid records are dropped, and every flow is projected onto the manifold \emph{before training and testing} (Online Resource~1, Section~S12); this pre-projection clamps byte totals above 1500 times the packet count, durations above 110~s and UDP header sizes, so the learned models never see such values. Records with identical features and label are then collapsed so that no flow appears on both sides of the split~\cite{engelen2021troubleshooting, lanvin2023errors}, which removes 39.1\%, 20.5\% and 42.7\% of the valid records and leaves 4,968,389, 1,658,869 and 123,508 flows (Online Resource~1, Table~S20). Records with identical features but different labels (149,526, 22 and 1,040) are kept as real ambiguity of the schema. The split into training, validation and test partitions is 70\%/10\%/20\%, stratified by label and family, random, and not grouped by time or host, so near-duplicates may fall on both sides and accuracies are optimistic.

The targets are XGBoost (CUDA), LightGBM (each 60 trees of depth 6, learning rate 0.1) and CatBoost (60 oblivious trees of depth 6). \emph{FFE-v2} is score-only; its snapping variant also reads thresholds. \emph{FFE-v1} (Online Resource~1, Algorithm~S1) is the reference baseline and reads thresholds for its snapping step. The \emph{projected white-box control} reads all trees: each step moves one feature across one split threshold, choosing the move that lowers the exact ensemble margin most, with every candidate projected onto the manifold before scoring; it counts two queries (the tree reads are free) and is cost-blind. For XGBoost it is a greedy leaf-tuple-style attack in the spirit of Zhang et al.~\cite{zhang2020efficient}; for LightGBM and CatBoost it is our generic version on the raw ensemble margin. It is not LT-Attack (whose official code needs a Bazel/C++ build we did not set up) and we did not run Cascioli et al.~\cite{cascioli2024faster}. On XGBoost we also compare with seven earlier baselines: a naive coordinate descent, \emph{ZOO-NIDS}~\cite{chen2017zoo}, \emph{HopSkipJump-NIDS}~\cite{chen2020hopskipjump}, \emph{MOEVA}~\cite{simonetto2024constrained, ghamizi2020moeva}, \emph{PSO}~\cite{han2021traffic} with our projection, and our re-implementations of \emph{CAPGD}~\cite{simonetto2024constrained} and a \emph{NIDSGAN-style generator}~\cite{zolbayar2022nidsgan}, which fit a small MLP surrogate on the training flows (surrogate computations are not counted). CAPGD was run on XGBoost in the single-split test and, in the seed study, also on LightGBM and CatBoost (three seeds); baselines use default budgets.

\subsection{Protocol, Metrics and Cohorts}
\label{sec:protocol}

All design choices and hyper-parameters of FFE-v2 were tuned (tuning queries are not counted) on cohorts of up to 1,000 detected malicious flows from the 10\% \emph{validation} split; the configuration was frozen and evaluated once on the benchmark's 1,000-flow \emph{test} cohorts (the identical flows used for FFE-v1 and the baselines, verified by the runner). This is not a clean hold-out: the FFE-v1 failures that motivated FFE-v2 were first observed on these test cohorts, re-segmentation was motivated by the control's moves on the same corpora, and the window pool uses the attacker's benign training sample. The five-seed study (retrained models, fresh cohorts, no tuning) is the cleaner confirmation. The cohorts are drawn once, with a fixed seed, from the malicious test flows that both XGBoost and LightGBM detect when clean (77.6\%, 99.8\% and 87.0\% of the test attacks of CSE-CIC-IDS2018, CIC-IDS2017 and UNSW-NB15), so all results are conditional on detection. For CatBoost the cohort is restricted to the flows CatBoost itself detects ($N=1000$, 997 and 983), so every success is a genuine evasion of that model. MOEVA and PSO need 183--445 queries per flow and attack the first 200 flows. FFE-v1 on XGBoost and LightGBM and the baselines are the saved benchmark outcomes; FFE-v1 on CatBoost and the control on LightGBM and CatBoost were run for this study.

ASR is the fraction of attacks with $P(\text{Malicious}\mid x^*)<0.5$; PFR the fraction of evasive candidates that satisfy the schema invariants; Real ASR counts an attack only when it both evades and is feasible. Feasibility is the schema verifier of Section~\ref{sec:formulation} (immutability, monotonicity, MTU and header bounds, 5\% tolerance on the coupled rates, identical for every method), which holds for FFE by construction; it does not check the wire-byte gap $x_4\ge x_{11}+20x_2$, which re-segmentation would violate if $x_4$ counted headers. ``100\% schema-feasible'' is relative to this verifier. We report Wilson 95\% intervals and exact McNemar tests on paired outcomes; no multiplicity correction is applied except where stated (nine cells, so individual $p$-values should be read with that in mind), and pooled seed tests treat flows as independent. Cost is the median over a method's own feasible evasions of the operational cost of Section~\ref{sec:cost} and of the normalized distance $D_{L_1}$; these medians cover different flows, so we add paired medians on flows two methods both evade. For the multi-seed study every seed retrains the targets on a new stratified split and redraws a cohort of up to 300 flows detected by XGBoost (LightGBM and CatBoost are attacked on the part of that cohort they also detect); the frozen configuration is used without tuning, with five seeds for XGBoost and three for LightGBM and CatBoost.
